\begin{proposicion}[Orden de los parámetros y de los errores]\label{prop:cortes-uniformes}
Con la banda histórica y los parámetros que siguen existe una constante \(C_\delta\)
común a todos los cortes admisibles. A escala \(n\), los errores de los apéndices
\ref{app:signed}, \ref{app:arrival}, \ref{app:interior} y \ref{app:weights} tienen
la siguiente descomposición. Escribimos \(o_b(1)\) para una cantidad que tiende a cero
cuando \(n\to\infty\) para cada \(b\) fijo.
\begin{enumerate}
\item En la identidad con signo y en el resto de llegada, el reemplazo de \(g^\flat\)
y de su exponencial cuesta a lo sumo \(C_\delta b+o_b(1)\).
Las salidas, la sustitución por menores aproximados, la cola de activaciones y los
restos de Taylor de la identidad detenida son \(o_b(1)\).
\item En la media interior son \(o_b(1)\) el error de inserción sumado, el cambio
de ley de los incrementos, la sustitución de los promedios uniformes de recorridos
y su suma de Riemann en los niveles.
\item En la inserción de un par, los restos de Taylor de la acción y del prefactor
contribuyen \(o_b(1)\) después de sumar en los pares. La sustitución de sus coeficientes
por las ocho integrales de estado a ese mismo corte también cuesta \(o_b(1)\).
\end{enumerate}
\end{proposicion}
\begin{proof}
Se fija la banda de~\cite[§5.4]{tc}, su \(\delta\), y se reservan los márgenes comunes
de difusión, vacancia, brecha, carga y salida \(c_*>2\beta\). La transición de clase
fija \(C_{\rm pres}\) antes del corte. Con esos márgenes se fijan los radios de los
menores y de sus campos, el umbral \(\eta\) y la base \(C_1\) del recuento tanto
determinista como integrado hasta \(8\log n\). La reserva exponencial para prescripciones
largas precede a la fracción de truncación del Taylor y al umbral \(\eta\).
Elíjase \(0<a<1\) con \(a\log(\max(2,C_1))<1/4\); después \(\rho\ge2\) con
\(a\log\rho>3\); por último \(\alpha<\alpha_{\rm hist}\) con \(2C\alpha<\eta/(2\rho^2)\)
y con la holgura de la clase exterior y del balanceo de torres.
Elijamos \(\delta_*<1\) con \(C_1\delta^{-2}\delta_*^{1/3}<1\) y
\(C_2\delta^{-2}\delta_*<1\) para las dos series de~\eqref{eq:llegada-nu}.
Ahora puede elegirse cualquier \(0<b<\min(\delta/2,\delta_*)\) suficientemente pequeño
para que \(x_0<1\) y se cumpla la desigualdad exponencial del
Lema~\ref{lem:llegada-integrada}. La prueba del Lema~\ref{lem:prescripcion-larga}
muestra dónde se paga con ella la parte positiva \(\beta j\) del peso de una línea larga.
Finalmente \(n_0(b)\) asegura \(j\ge\sqrt n,\ v^2/j\le\delta_*\),
\(C(4v+12)\le j^{1/3}\) hasta \(v=8\log n\), y todas las cotas de clase en el hijo.
Los radios globales no aumentan tras fijar la carga.

Para el primer punto, las dos series de tipos son geométricas una vez extraído
\(b/\delta_*\). Sus bases y el multiplicador acotado de la transferencia a \(\mu\)
dependen sólo de las elecciones anteriores a \(b\), aumentando \(n_0(b)\) para que
\(e^{|A-R_Q(n)|}\le2\). El coste restante es
\[
 O_b\!\left(nj^{-8}+n^{-1}+n(1+j)e^{-2j}+ne^{-c_b n}
           +n e^{-c_\delta n+D_\delta n^{2/3}}\right)=o_b(1).
\]
Los restos cuadráticos en la identidad con signo son \(O_b(n^{-1})\) a escala \(n\).
En el segundo punto, los residuos cúbicos cuestan \(O_b(n\sum_{m>j}m^{-3})\);
los incrementos centrados y el cambio de ley cuestan
\(O_b(n(\sum_{m>j}m^{-4})^{1/2})\). Ambos tienden a cero. Si la aproximación
uniforme de la media tiene error \(\epsilon_n(b)\to0\), su suma cuesta
\(n\epsilon_n(b)\sum_{m>j}m^{-2}=o_b(1)\). Para obtener esa aproximación se fija
primero la longitud de los recorridos, se toma \(n\to\infty\) y después se usa la
cola geométrica de la resolvente en \([b,1]\). La continuidad sobre ese intervalo
hace tender a cero la diferencia de Riemann. En el tercer punto el error puntual en
el cociente de inserción es \(O_b(j^{-3})\); cada par lleva masa \(O_b(j^{-2})\)
y hay \(O(j^3)\) pares. El coste total es \(O_b(n/j^2)=o_b(1)\).
Las coincidencias accidentales de los recorridos truncados cuestan \(O_{b,N}(1/n)\)
y los términos más largos están dominados por \(C_b\sum_{q>N}(1+q)^2r_b^{2q-O(1)}\),
\(r_b<1\). Estos límites justifican las sustituciones a corte fijo.
\end{proof}
Al comparar después la frontera a corte \(b\) con su valor final aparece
\(\eta(b)=\limsup_n|\mathcal M_n(b)-(P_T+J-S)|\to0\) al bajar \(b\),
como se usa en el apéndice~\ref{app:closing}. Este último error pertenece al límite
del perfil de frontera en \(b\); la proposición sólo emplea \(o_b(1)\) para el límite en \(n\).
